\documentclass[10pt]{article}

\usepackage[utf8]{inputenc}

\usepackage[T1]{fontenc}

\usepackage{amsmath}

\usepackage{amsfonts}

\usepackage{amssymb}

\usepackage[version=4]{mhchem}

\usepackage{stmaryrd}



\begin{document}

Всякий интеграл $\int R(x, y) d x$, где $y-$ функция от $x$, определяемая уравнением $F(x, y)=0$, задающая алгебраич. У. к., а $R(x, y)$ - рациональная функция, приводится к ивтегралу от рациональных функций и выражается в элементарных функциях. М. И. Войцеховский.



УНИМОДАЛЬНОЕ РАСПРЕДЕЛЕНИЕ, о Д н ов е р п̈ и но е р а с II р е д е л е ни е,- вероятностная мера на прямой, функция распределения к-рой $F(x)$ выпукла при $x<a$ и вогнута при $x>a$ для нек-рого действительного $a$. Число $a$ в этом случае наз. модой (в е р ш и о й) и определяется, вообще говоря, неоднозначно; точнее, множество мод (вершин) данного У. р. образует замкнутый интервал, возможно, вырожденный.



Примерами У. р. служат нормальное распределение, равномерное распределение, Коши распределение, Стьюдента распределение, «Хи-квадрат» распределение. А. Я. Хинчин [1] получил следуюший к р и т е р и й у н и м о д а л ь н о с т и: для того чтобы функция $f(t)$ была характеристич. функциеї У. р. с модой в нуле, необходимо и достаточно, чтобы она допускала представление



\[

f(t)=\frac{1}{t} \int_{0}^{t} \varphi(u) d u, \quad f(0)=1,

\]



где $\varphi(u)$ - нек-рая характеристич. функция. В терминах функций распределения это равенство эквивалентно следующему



\[

F(x)=\int_{0}^{1} G\left(\frac{x}{u}\right) d u

\]



где $F(x)$ и $G(x)$ соответствуют $f(t)$ и $\varphi(t)$. Иными словами, $F(x)$ ун и м о д а л ь н а тогда и только тогда, когда она является функцией распределения произведения двух независимых случайных величин, одна из к-рых распределена равномерно на отрезке $[0,1]$.



Для распределений, к-рые задаются своими характеристич. функциями (как, напр., для устойчивого распределения), проверка факта их унимодальности представляет трудную аналитич. задачу. Кажущийся естественным путь представления данного распределения как предела композиции У. р. не приводит к цели, т. к., вообще говоря, композиция двух У. р. не будет У. р. (хотя для симметричных распределений унимодальность сохраняется гіи композиции и долгое время казалось, что так будет в общем случае). Напр., если $F$ имеет плотность



\[

p(x)=1, \quad 0<x \leqslant \frac{5}{6} \text { и } p(x)=0

\]



в др. случаях, то плотность композиции $F$ с $F$ имеет два максимума. Поэтому было введено (см. [2]) Іонятие сильной унимодальности (сильной одновершинности): распределение наз. с и л ь н о ун и м о д а л ь н ы м (с ильно о днов е р ши ины м), если его свертка с любым У. р. унимодальна. Всякое сильно унимодальное распределение является У. р.



Решетчатое распределение, приписывающее точкам $a+h k, k=0, \pm 1, \ldots ., h>0$, вероятности $p_{k}$ наз. У. р., если существует такое целое $k_{0}$, что $p_{k}$ как функция от $k$ является неубывающей при $k \leqslant k_{0}$ и невозрастающей при $k \geqslant k_{0}$. П римерами решетчатых У. р. являются Пуассона распределение, биномиальное распределение, геометрическое распределение.



Нек-рые результаты, касающиеся распределений, могут быть усилены в лредположении унимодальности. Так, напр., Чебыиева неравенство для случайной величины $\xi$, имеющей У. р., можно уточнить следующим образом



\[

\mathrm{P}\left\{\left|\xi-x_{0}\right| \geqslant k \zeta\right\} \leqslant \frac{4}{9 k^{2}}

\]



для любого $k>0$, где $x_{0}-$ мода, а $\zeta^{2}=\mathrm{E}\left(\xi-x_{0}\right)^{2}$. JІит.: [1] Х и н ч и н А. Я., Об унимодальных распределениях, “Изв. НИИ матем и мех. ун-та", Томск, 1938, т. 2, в. 2,

с. $1-7$, [2] И б р а г и м о в И. А., О композиции одновершинных распределений, “Теор. вероятн. и ее примен.», 1956 , т. 1, в. 2, с. 283-88, [3] Ф е л л е р В., Введение в теорию вероятностей и ее цриложсния, пер. с англ., 2 изд., т. 2, М., 1967.



УНИМОДУЛЯ РНАЯ ГРУППА - топологическая группа, левоинвариантная $X$ аара мера на к-роӥ правоинвариантна или, что равносильно, инвариантна относительно преобразования $a \longmapsto a^{-1}$. Группа Ли $G$ унимодулярна тогда и только тогда, когда



$|\operatorname{det} \operatorname{Ad} g|=1 \quad(g \in G)$,



где Аd - присоединенное представление. Для связных груип Ли $G$ это равносильно тому, что tr ad $x=0$ $(x \in \mathfrak{q})$, где ad - присоединенное представление алгебры Ли g группы $G$. Любая компактная, дискретная или абелева локально компактная группа, а также любая связная редуктивная или нильпотентная группа Ли унимодулярна.



А. Л. Онищик.



УНИМОДУЛЯРНАЯ МАТРИЦА - квадратная матрица, определитель к-роӥ равен \textbackslash pm 1 . Иногда, при рассмотрении матриц над коммутативным кольцом, под У. м. понимают обратимую матрицу. O. А. Иванова.



УНИМОДУ ЛЯ ЯНОЕ ПРЕОБРАЗОВАНИЕ - лИНеӤное преобразование конечномерного векторного пространства, матрица к-рого имеет определитель \textbackslash pm 1 .



О. А. Иванова.



УНИПОТЕНТНАЯ ГРУППА - подгруппа $U$ линейной алгебраич. группы $G$, состоящая из унипотентных элементов. Если отождествить $G$ с ее образом при изоморфном вложении в группу $G L(V)$ автоморфизмов подходящего конечномерного векторного пространства $V$, то У. г. - это подгруппа, лежащая в множестве



\[

\left\{g \in G L(V) \mid(1-g)^{n}=0\right\}, \quad n=\operatorname{dim} V,

\]



всех унипотентных автоморфизмов пространства $V$. Зафиксировав в $V$ базис, можно отождествить $G L(V)$ с полной матричной группой $G L_{n}(K)$, где $K$ - основное алгебраически замкнутое поле; матричная группа $U$ также называется при этом У. г. Пример У. г.группа $U_{n}(K)$ всех верхнетреугольных матриц из $G L_{n}(K)$ с единицами на диагонали. Если $k$ - шодполе поля $K$ и $U$ - унипотентная подгруппа в $G L_{n}(k)$, то $U$ сопряжена над $k$ с нек-рой подгруппой группы $U_{n}(k)$. В частности, все элементы из $U$ имеют в $V$ общий ненулевой неподвижный вектор, а $U$ является нильпотентной группой. Эта теорема показывает, что алгебраические У. г.- это в точности замкнутые (в топологии Зариского) подгрупшы групп $U_{n}(K)$ при различных $n$.



В любой линейной алгебраич. группе $H$ имеется единственная связная нормальная унипотентная подгрупіа $R_{u t}(H)$ (ун и п о т е н т ны й р а д и к а л), фактор $H / \stackrel{R}{R}_{n}(H)$ по к-рому редуктивен. Это в известной степени сводит изучение строения любой групшы к изучению строения редуктивных групп и У. г. В отличие от редуктивного случая классификация алгебраических У. г. к настоящему времени (1984) неизвестна.



Всякая подгрупша и всякая факторгруппа алгебраической У. г. снова являіотся У. г. Если char $K=0$, то $U$ всегда связна, причем әкспоненциальное отображение ехр $: \breve{u} \rightarrow U$ (где $\breve{u}$ - алгебра Ли группы $U$ ) является изоморфизмом алгебраич. многообразий; если же char $K=p>0$, то существуют несвязные алгебраические У. г.: напр., аддитивная группа $\mathbf{G}_{a}$ основного поля $K$ (ее можно отождествить с $U_{2}(K)$ ) является $p$-группой и потому содержит конечные У. г. В связной У. г. $U$ имеется такая последовательность нормальных делителей $U=U_{1} \supset U_{2} \supset$. . $\supset U_{s}=\{e\}$, что все факторы $U_{i} / U_{i+1}$ одномерны. Всякая связная алгебраическая одномерная У. г. изоморфна $\mathbf{G}_{a}$. Это сводит изучение связных алгебраических У. г. к описанию кратных расширений групп типа $\mathbf{G}_{a}$.





\end{document}